\section{Produção: robustez, informação e calibração condicional}
\label{sec:c09-producao}

A produção usa o experimento periódico de Fourier definido anteriormente,
com parâmetros auxiliares conhecidos. Foram geradas e verificadas
\(5\,396\) observações, analisadas em \(25\,956\) posteriores de base.
O painel pareado de prioris contém \(1\,600\) integrações sobre os mesmos
\(32\) dados centrais, dez análises e cinco medidas: \(320\) integrações
reavaliam a priori uniforme já existente e \(1\,280\) são contrastes
adicionais. Não são novas observações nem uma calibração SBC das prioris
reponderadas. As sementes, os identificadores, as configurações e os
resultados preservam também os casos numericamente indeterminados.

As tabelas positivas interpolam log-verossimilhança em
\(\alpha=\arcsin u\), por PCHIP, e integram dentro de cada intervalo por
quadratura positiva. Malhas aninhadas, ordens 32 e 16 e referências
independentes selecionadas verificam os funcionais. A concordância nesses
testes tem alcance finito: não constitui limite uniforme de erro da
resposta física. Na síntese dos quantis, uma análise sem aprovação dos
critérios de malha e resposta recebe \([0,1]\); nenhum ID é eliminado.
Os intervalos descritivos deste capítulo propagam a incerteza numérica,
não a incerteza amostral de uma média populacional.

\subsection{Sensibilidade à medida e ao suporte}

No painel de mesmos dados, a mediana de \(Q_{0,95}\) de A0 é
\([0,9355;0,9360]\) para a priori uniforme em \(u\),
\([0,9670;0,9675]\) para a uniforme em \(u^2\) e
\([0,6856;0,6861]\) para a logarítmica com corte \(10^{-3}\).
Alterar apenas esse corte para \(10^{-4}\) ou \(10^{-2}\) fornece,
respectivamente, \([0,6184;0,6189]\) e \([0,7663;0,7668]\).
As medianas de KL correspondentes ficam entre \(0,0261\) e
\(0,0489\) nat. A proximidade de alguns quantis aos quantis prévios
é acompanhada por medidas de informação; não é tomada isoladamente como
prova de ausência de aprendizado.

Para B construído das observações CN, com covariância completa e variável,
a priori uniforme produz mediana de \(Q_{0,95}\) em
\([0,9473;0,9478]\) e KL mediana de \(0,00230\) nat. Com corte
logarítmico \(10^{-4}\), os valores são \([0,6165;0,6170]\) e
\(0,00166\) nat. Assim, neste experimento condicional e neste painel,
a escolha de medida e suporte altera o limite muito mais que a pequena
distância informacional típica da posterior comprimida à sua própria
priori. Isso não estabelece um teorema de ausência de informação em PTAs.
W1, diferenças de CDF e razões de odds para \(u\leq0,2\) estão também
registrados por dado; a interpretação não depende de um único quantil.

\subsection{Fatorial da covariância}

O desenho separa covariância completa/diagonal de fixa/variável e conserva
o termo log-determinante da densidade normalizada. São \(24\) contrastes
de \(500\) pares, cobrindo três prioris e as duas origens de observações B.
Para B de origem CN e priori uniforme, diagonal menos completa, ambas
variáveis, tem diferença média de \(Q_{0,95}\) em
\([-0,00047;0,00051]\). Há \(17\) pares com diferença negativa,
\(14\) positiva e \(469\) com sinal não resolvido numericamente.
Com priori logarítmica, esses números são \(175\), \(174\) e \(151\).
A diagonalização pode produzir efeitos de ambos os sinais; não é
universalmente conservadora. Médias compatíveis com zero na resolução
usada não demonstram equivalência entre os procedimentos.

A escolha fixa/variável é um eixo distinto. No mesmo B de origem CN,
com priori uniforme e matriz completa, fixa menos variável apresenta
média em \([0,00020;0,00118]\). Em B de origem gaussiana com priori
logarítmica, a média correspondente é negativa,
\([-0,00133;-0,00036]\). São descrições das realizações registradas,
sem intervalo de confiança populacional ou extrapolação a outros sinais.

\subsection{Contaminantes: limites mais baixos podem ser incorretos}

Quatro cenários usam \(64\) realizações em \(u_*=0,5\), com
contaminante monopolar ou dipolar de amplitude relativa \(0,25\) ou \(1\).
A comparação pareada conserva a observação e omite ou inclui o componente
conhecido na covariância. A Tabela~\ref{tab:c09-contaminantes-producao}
apresenta os contrastes de A0; os \(12\) contrastes completos abrangem
também as duas análises B.

\begin{table}[htbp]
\centering
\caption{Média descritiva de \(\Delta Q_{0,95}\), omissão menos inclusão,
em A0. Intervalos numéricos; \(64\) pares por linha.}
\label{tab:c09-contaminantes-producao}
\begin{tabular}{lcc}
\hline
Componente & Amplitude relativa & Intervalo da média\\
\hline
Monopolar & \(0,25\) & \([-0,06715;-0,06617]\)\\
Monopolar & \(1\) & \([-0,18319;-0,18221]\)\\
Dipolar & \(0,25\) & \([-0,04852;-0,04754]\)\\
Dipolar & \(1\) & \([-0,30819;-0,30721]\)\\
\hline
\end{tabular}
\end{table}

No cenário dipolar de amplitude relativa unitária, A0 apresenta mediana
de quantil \([0,9504;0,9509]\) com inclusão e
\([0,6655;0,6660]\) com omissão. Em \(61\) dos \(64\) pares o
quantil diminui com sinal resolvido. Um limite aparentemente mais forte
pode, portanto, decorrer de especificação incorreta. Em B de origem CN,
a diferença média nesse cenário é \([-0,01264;-0,01167]\): a resposta
menor à omissão também não demonstra robustez observacional, especialmente
quando a compressão já retém pouca informação sobre a massa.

\subsection{Ruído, espectro e variabilidade entre realizações}

Os dois cenários de ruído contêm \(500\) realizações cada. Em A0, as
medianas de KL são \(0,01121\) nat sob ruído vermelho forte e
\(0,01920\) nat sob ruído branco forte. Em B de origem CN são
\(0,00284\) e \(0,00042\) nat. As medianas dos quantis ficam próximas
de \(0,95\), mas os percentis entre realizações e as KL distinguem
as distribuições. As realizações de cenários diferentes têm sementes
próprias; diferenças entre essas medianas não são contrastes pareados.

Nos cenários espectrais, com \(64\) realizações por combinação de inclinação
e massa fixa, o expoente \(3\) fornece KL mediana de A0 de
\(0,00752\) nat em \(u_*=0,5\). Com expoente \(5,5\), ela passa a
\(0,09961\) nat; a mediana de \(Q_{0,95}\) passa de aproximadamente
\(0,9502\) a \(0,9141\). Perto do corte, \(u_*=0,995\), o espectro
\(5,5\) fornece KL mediana \(0,26076\) nat e quantil mediano
\([0,9751;0,9756]\). Não se aplica a hipótese uniforme de SBC a esses
cenários de massa fixa. Os \(128\) dados por massa do painel de fronteiras
completam a descrição da variabilidade condicional.

\subsection{Distâncias e indeterminação numérica}

O cenário de distâncias sorteia uma escala global entre \(0,9\), \(1\)
e \(1,1\), com pesos \(0,25\), \(0,50\) e \(0,25\). A inferência
correta soma as densidades completas, após todos os canais; não mistura
covariâncias nem sorteia uma escala independente por frequência. Compara-se
essa mistura à análise com escala nominal sobre os mesmos \(500\) dados.

As \(3\,000\) posteriores foram integradas, mas \(834\) análises A0
falharam no controle pontual de interpolação da resposta em logL.
Preservando esses casos como quantis \([0,1]\), o contraste médio
nominal menos mistura em A0 fica em \([-0,90362;0,76454]\), inconclusivo.
Nas duas análises B, os \(500\) contrastes têm sinal não resolvido e
média em cerca de \([-0,00049;0,00049]\). A resolução adotada não permite
afirmar equivalência, nem que incertezas de distância sejam irrelevantes.

\subsection{Calibração da massa e dos eventos de likelihood}

A família prospectiva contém \(18\) grupos de \(500\) observações e
\(126\) testes. Foram conservados todos os IDs, inclusive os intervalos
indeterminados \([0,1]\). No conjunto de \(9\,000\) análises SBC,
\(8\,621\) PIT de massa e \(8\,179\) PIT de logL têm resolução
operacional. Os eventos usam o limiar avaliado diretamente na resposta
física da massa verdadeira. As regiões desconectadas são integradas na
tabela positiva; patamares seriam tratados separadamente como átomos.
Nenhum átomo foi encontrado nos interpolantes desta produção.

A troca do limiar interpolado pelo físico é relevante: uma diferença de
logL de apenas \(1,2091\times10^{-5}\) alterou uma probabilidade de evento
em \(0,139777\). Assim, uma tolerância pontual de logL não é, sozinha,
um limite de erro para a probabilidade do evento. Após os controles de
discretização, de resposta e de limiar, a família completa com correção de
Holm a \(0,05\) apresenta \(116\) não rejeições para todos os valores
admissíveis dos intervalos fornecidos e \(10\) testes numericamente
indeterminados. Não há rejeição robusta nesta síntese condicional.
Os envelopes observados não possuem certificado físico uniforme:
não rejeitar não prova calibração física completa, aprendizado da massa
ou equivalência entre modelos. Os registros da síntese preliminar sem
eventos permanecem preservados como histórico.

\subsection{Escala física e alcance das conclusões}

Como \(m_gc^2=(h/T)u\), todo quantil relatado neste experimento é
\(Q_{0,95}(m_gc^2)=(h/T)Q_{0,95}(u)\). Sob priori uniforme e likelihood
constante, ele já vale \(0,95h/T\). Portanto, proximidade de um limite
à escala \(h/T\), como no benchmark MeerKAT discutido na revisão, precisa
ser interpretada com a medida, o suporte, as frequências e o modelo
probabilístico completos. Os resultados do painel de prioris ilustram
quantitativamente esse problema, sem reanalisar as observações MeerKAT.

Os dados sintéticos periódicos não incorporam automaticamente cadência
irregular, ajuste do modelo de temporização, acoplamento entre frequências
ou possível impropriedade complexa dos coeficientes obtidos de TOAs reais.
Também não marginalizam os parâmetros auxiliares aqui conhecidos.
Comparações de evidência entre espaços de observação diferentes não são
interpretadas como fatores de Bayes. A extensão escalar e a aproximação
observacional permanecem tarefas próprias de C10 e C11.

A consolidação está em \path{results/C09/robustness_synthesis/}, com
\(90\) grupos descritivos, \(39\) contrastes e \(14\,268\) diferenças
pareadas. O acumulado C09 é de \(67\,072\,009\) avaliações de likelihood;
o acumulado físico D3 é de \(7\,920\) nós e aproximadamente
\(85,696\) trilhões de produtos reais estimados. Falhas e tentativas
interrompidas permanecem contabilizadas. Foram restaurados e conferidos
\(24\,238\) arquivos lógicos a partir dos pacotes da produção e de suas
dependências. O manifesto da versão vincula fontes, resultados e PDF.
